%
\documentclass[12pt,a4letter]{article}
%Các gói lệnh cơ bản
\usepackage{amsmath,amsxtra,amssymb,latexsym, amscd,amsthm}
\usepackage{indentfirst}
\usepackage[mathscr]{eucal}
\usepackage{graphicx}
\usepackage{makecell}
\usepackage{longtable}
\usepackage[utf8]{vietnam}
%Kích thước của trang đề thi
\voffset=-2cm
\hoffset=-2cm
\textheight 23truecm 
\textwidth 17truecm 
%Dùng gói lệnh cho câu hỏi và lời giải
% \usepackage[nosolutionfiles]{answers}
\usepackage{answers}
\theoremstyle{definition}
\newtheorem{cauhoi}{Câu}
\Newassociation{traloi}{Soln}{goiytraloi}
\renewcommand{\Solnlabel}[1]{{\bf Lời giải #1}.}

\newenvironment{phandiem}{%
\par\noindent
\begin{longtable}[t]{|p{15cm}|p{1.2cm}|}
\hline
}{%
\end{longtable}
}
\newcommand\diem[1]{ & \hfil \textbf{#1} \\ \hline}
\newcommand\tongdiem[1]{\hfill [\textbf{#1} điểm]}
%Định nghĩa tiêu đề
\def\tendhqg{\small \textbf{ĐẠI HỌC QUỐC GIA HÀ NỘI}}
\def\tentruong{\small \textbf{ĐẠI HỌC KHOA HỌC TỰ NHIÊN}}
\def\tenkythi{ĐỀ THI KẾT THÚC MÔN HỌC}
\def\namhoc{NĂM HỌC 2017-2018}
\def\tenmonhoc{Tối ưu hóa}
\def\mamonhoc{MAT2407}
\def\sotinchi{3}
\def\svkhoa{60MTTT}
\def\nganhhoc{MT\& KHMT, Toán - Tin ứng dụng}
\def\thoigian{90 phút }
\def\deso{1}
\def\dapan{ĐÁP ÁN VÀ THANG ĐIỂM}
\def\nguoilamdapan{PGS.TS. Nguyễn Hữu Điển}
\def\ngaylam{Hà nội, ngày 10 tháng 5 năm 2018}

\usepackage{mathpazo}
\usepackage{color}
\graphicspath{{hinh/}}
\begin{document}
\setlength{\baselineskip}{16truept}
%Mở tệp ghi đáp án
\Opensolutionfile{goiytraloi}[traloi] 
%Phần tiêu đề của đề thi
\thispagestyle{empty}
\begin{minipage}[b]{0.4\textwidth}
\centering
{\bf \tendhqg}\\
{\bf \tentruong}\\ 
-------------\\
\quad \\
\end{minipage}
\hfill
\begin{minipage}[b]{0.6\textwidth}
\centering
{\bf \tenkythi}\\
{\bf  \namhoc}\\
------oOo-------\\ 
\qquad 
\end{minipage}

\centerline{\textbf{Môn thi:} \textbf{\tenmonhoc}}
\leftline{Mã môn học: \textbf{\mamonhoc}\ \hfil Số tín chỉ:  \textbf{\sotinchi} \hfil Đề số: \textbf{\deso} \hfil }
\leftline{Dành cho sinh viên khoá: \textbf{\svkhoa} \hfil Ngành học: \textbf{\nganhhoc}\hfil }
\centerline{Thời gian làm bài \textbf{\thoigian} (không kể thời gian phát đề)}
\noindent \hrulefill

%%%%%%%%%%%%%%%%%%%%%%%%%%%%%%%%%%%%
%Phần câu hỏi và trả lời
\begin{cauhoi}%1
(2 điểm) Một công ty may cần sản xuất 2000 quần và ít nhất 1000 áo để giao cho khách hàng. Để sản xuất số quần áo đó, công ty sử dụng những tấm vải giống nhau và mỗi tấm vải có 6 cách cắt ra số quần, áo cho trong bảng sau
\begin{center}
\begin{tabular}{| c | c  c  c  c  c  c|}
\hline
Cách cắt&1&2&3&4&5&6\\
\hline
Số quần&90&80&70&60&120&0\\
\hline
Số áo&35&55&70&90&0&100\\
\hline
\end{tabular}
\end{center}

1. Hãy lập mô hình toán học của bài toán 
xác định số mảnh vải phải cắt theo mỗi cách sao cho số quần áo đủ để giao cho khách hàng và tổng số tấm vải phải sử dụng ít nhất.

2. Đưa mô hình toán học trên về bài toán đối ngẫu.

%%%%%%%%%%%%%%%Lời giải
\begin{traloi}\tongdiem{2}
\begin{phandiem}
Gọi $x_1, x_2, x_3, x_4, x_5, x_6$ lần lượt là số vải phải cắt theo cách   $1, 2, 3, 4, 5, 6$ để thỏa mãn các yêu cầu của bài toán.

Ta có điều kiện cho ẩn số là  $x_j\ge 0, j=1, 2, 3, 4, 5, 6$. 

Tổng số quần áo thu được là :
 
Quần :  $90x_1+80x_2+70x_3+60x_4+120x_5$

Áo : $35x_1+55x_2+70x_3+90x_4+100x_6$

Để số quần áo sẽ sản xuất đủ để giao cho khách hàng thì ta phải có điều kiện

  $90x_1+80x_2+70x_3+60x_4+120x_5=2000$

 $35x_1+55x_2+70x_3+90x_4+100x_6\ge 1000$.

\diem{1}

Tổng số tấm vải phải sử dụng để sản xuất ra số quần áo trên là

 $x_1+x_2+x_3+x_4+x_5+x_6$ 

Để tổng số tấm vải phải sử dụng để sản xuất ra số quần áo trên ít nhất thì ta có điều kiện

 $$x_1+x_2+x_3+x_4+x_5+x_6\rightarrow\min$$

Như vậy mô hình toán học của bài toán trên là bài toán QHTT sau
\begin{center}
$ F(x)=x_1+x_2+x_3+x_4+x_5+x_6\rightarrow\min$ \\
 $ \begin{cases}
  90x_1+80x_2+70x_3+60x_4+120x_5&=2000\\
 35x_1+55x_2+70x_3+90x_4+100x_6&\ge 1000\\
x_j\ge 0, j=1, 2, 3, 4, 5, 6.&\\ 
\end{cases} $
\end{center}
\diem{1}

\end{phandiem}
\end{traloi}
\end{cauhoi}


\begin{cauhoi}%2
(4 điểm)  
Cho bài toán quy hoạch tuyến tính
\begin{center}
$ F(x)=27x_1+50x_2+18x_3\rightarrow\max $ \\
 $ \begin{cases}
x_1+2x_2+x_3&\le 2\\ 
-2x_1+x_2-x_3&\le 4\\ 
x_1+2x_2-x_3&\le -2\\ 
x_3\le 0.&\\ 
\end{cases} $
\end{center}

i) Thiết lập bài toán đối ngẫu với bài toán trên;

ii) Giải bài toán đối ngẫu vừa thiết lập bằng phương pháp đơn hình mở rộng; Rút ra kết luận nghiệm của bài toán đã cho.

%%%%%%%%%%%%%%%Lời giải
\begin{traloi}\tongdiem{4}
\begin{phandiem}
i) Chuyển về dạng bài toán đối ngẫu:  
\begin{center}
$L(y)=2y_1+4y_2-2y_3\rightarrow\min$\\
$\begin{cases}
y_1-2y_2+y_3&= 27\\ 
2y_1+y_2+2y_3&=50\\ 
y_1-y_2+2y_3&\le 18\\ 
y_i\ge 0 (i=&1,2,3)\\ 
\end{cases} $
\end{center}
\diem{1}

ii) Giải bằng phương pháp đơn hình bài toán trên: Đưa thêm 
	$y_4$ biến phụ và $y_5, y_6$ là biến giả, bài toán trở thành

\begin{center}
$L(y)=2y_1+4y_2-2y_3+My_5+My_6\rightarrow\min$\\
$\begin{cases}
y_1-2y_2+y_3+y_5&= 27\\ 
2y_1+y_2+2y_3+y_6&=50\\ 
y_1-y_2+2y_3+y_4&= 18\\ 
y_i\ge 0 (i=1,2,3,4,5,6)&\\ 
\end{cases} $
\end{center}
\diem{1}

\begin{center}
\renewcommand{\arraystretch}{1.25}
\begin{tabular}{| l | l | l | l  l  l  l  l  l |}
\hline
$J$&$c_J$&$x_J$&1&2&3&4&5&6\\ 
\cline{4-9}
   &     &   &2&4 &-2 &0&$M$&$M$\\ 
\hline
5&$M$&27&1 &-2 & 1   &0& 1&0\\
6&$M$&50 &2& 1&[2]  & 0& 0&1\\
4&0     &18 &1& 1& -1   & 1& 0&0\\
\hline
&L(x) &0  &-2&-4&2& 0& 0&0\\
&       &77 & 3&-1&[3]&0& 0&0\\
\hline
\end{tabular}
\end{center}
\diem{1}

\begin{center}
\renewcommand{\arraystretch}{1.25}
\begin{tabular}{| l | l | l | l  l  l  l  l  l |}
\hline
5&$M$&2   &0&-5/2 &0   &0& 1&\\
3&-2    &25 &1& 1/2 &1  & 0& 0&\\
4&0     &43 &2 &-1/2& 0   & 1& 0&\\
\hline
&L(x) &-50  &-4&-5&0& 0& 0&\\
&        &2    & 0&-5/2&0&0& 0&\\
\hline
\end{tabular}
\end{center}

Trong bảng đơn hình thứ hai ta có $\bar c_j\le 0$ với mọi $j$ nên bài toán mở rộng có phương án tối ưu là; $y^0=(0,0,25,43,2,0)$.

Do phương án của bài toán tối ưu có $y_5=2>0$ nên bài toán đối ngẫu không có phương án tối ưu, từ đây suy ra bài toán gốc cũng không có phương án tối ưu. 
\diem{1}

\end{phandiem}
 \end{traloi}
\end{cauhoi}

\Writetofile{goiytraloi}{\protect\newpage}

\begin{cauhoi}%3
(4 điểm) Cho số liệu bài toán vận tải và ma trận cước phí $C$ như sau:

$a=(80, 70, 100)^T$ \qquad $b=(40, 100, 60, 50)^T$

$C=\begin{bmatrix}
1&2&4&3\\ 
2&4&5&1\\ 
4&1&2&5\\ 
\end{bmatrix}$ 

a) Lập bảng vận tải cho bài toán trên.

b) Tìm phương án xuất phát theo phương pháp cước phí nhỏ nhất.

c) Giải bài toán vận tải trên theo phương pháp quy không cước phí.

%%%%%%%%%%%%%%%Lời giải
\begin{traloi}\tongdiem{4}
\begin{phandiem}

\noindent 1. Đây là bài toán vận tải dạng $\min$. Nơi cung cấp hàng $a_1=80$, $a_2=70$, $a_3=100$,và nơi nhận hàng $b_1=40$, $b_2=100$, $b_3=60$, $b_4=50$. Dễ thấy đây là bài toán vận tải cân bằng thu phát.

2.  Phương án cơ sở theo phương pháp góc tây bắc trong  bảng ~1.

\begin{center}
\begin{tabular}{|c|ll|ll|ll|ll|}%
\hline
\diaghead{\theadfont Diag Col}%^^A
 {$A_i$}{$B_j$}&
$B_1:$&40&$B_2:$&100&$B_3:$&60&$B_4:$&50\\
\hline
                 &1 &    & 2&    &4&  &   3 &  \\
 $A_1:80$  &X  &40&   X  &0& X  &40&   &\\
\hline
                 &2 &    & 4&    &5&  &   1 &  \\
 $A_2:70$  &  &   &    &&  X &20&X   &50\\
\hline
                 &4 &    & 1&    &2&  &   5 &  \\
 $A_3:100$  &  &   &   X &100&   &   &   &\\
\hline
\end{tabular} \\
Bảng 1. Phương án cơ sở 
\end{center}

\diem{1}

\noindent  2. Vòng lặp thứ nhất:

{\it Bước 1.} Phương án $x$ không thoái hóa vì $|G|=m+n-1=6$. Bổ sung ô $(1,2)$

Tính $\bar c_{ij}$ như bảng~2.


\begin{center}
\begin{tabular}{|c|ll|ll|ll|ll|l|}%
\hline
\diaghead{\theadfont Diag Col}%^^A
 {$A_i$}{$B_j$}&
$B_1:$&40&$B_2:$&100&$B_3:$&60&$B_4:$&50&\\
\hline
                 &0 &    & 0&   + &0&  -&   3 &  &$r_1=0$\\
 $A_1:80$  & X &60&  X   &0& X  &40&   &&\\
\hline
                 &0 &    & 1&    &0&  &   0 &&$r_2=1$  \\
 $A_2:70$  &  &    &    &  &  X &20&  X &50&\\
\hline
                 &4 &    & 0&  - &-1&  +&   6 & &$r_3=-1$ \\
 $A_3:100$  &  &   & X   &100&   &   &   &&\\
\hline
                 &  $s_1=$&1   & $s_2=$&2   & $s_3=$  &4&$s_4=$ &0&\\
\hline
\end{tabular} \\
{Bảng 2. Bước 1}
\end{center}

Ô vi phạm dấu hiệu tối ưu (3,3) là ô đưa vào.
\diem{1}

{\it Bước 2.} 

\begin{center}
\begin{tabular}{|c|ll|ll|ll|ll|l|}%
\hline
\diaghead{\theadfont Diag Col}%^^A
 {$A_i$}{$B_j$}&
$B_1:$&40&$B_2:$&100&$B_3:$&60&$B_4:$&50&\\
\hline
                 &0 & -   & 0& +   &1&  &   4 &  &$r_1=0$\\
 $A_1:80$  & X &40&  X   &40&   &&   &&\\
\hline
                 &-1 & +   & 0&    &0&-  &   0 &&$r_2=1$  \\
 $A_2:70$  &  &    &    &  &  X &20&  X &50&\\
\hline
                 &4 &    & 0& - &0& + &   7 & &$r_3=0$ \\
 $A_3:100$  &  &   & X   &60& X  &40   &   &&\\
\hline
                 &  $s_1=$&0   & $s_2=$&0   & $s_3=$  &-1&$s_4=$ &-1&\\
\hline
\end{tabular} \\
{Bảng 3. Bước 2}
\end{center}

\diem{1}

{\it Bước 3.} 

Tính $\bar c_{ij}$ như bảng 3.

Ô vi phạm dấu hiệu tối ưu (2,1) là ô đưa vào. 


\begin{center}
\begin{tabular}{|c|ll|ll|ll|ll|l|}%
\hline
\diaghead{\theadfont Diag Col}%^^A
 {$A_i$}{$B_j$}&
$B_1:$&40&$B_2:$&100&$B_3:$&60&$B_4:$&50&\\
\hline
                 &0 &    & 0&    &1&  &   3 &  &$r_1=0$\\
 $A_1:80$  & X &20&  X   &60&   &&   &&\\
\hline
                 &0 &    & 1&    &1&  &   0 &&$r_2=-1$  \\
 $A_2:70$  & X & 20   &    &  &   &&  X &50&\\
\hline
                 &4 &    & 0&  &0&  &   6 & &$r_3=0$ \\
 $A_3:100$  &  &   & X   &40& X  &60   &   &&\\
\hline
                 &  $s_1=$&0   & $s_2=$&0   & $s_3=$  &0&$s_4=$ &1&\\
\hline
\end{tabular} \\
{Bảng 4. Bước 3}
\end{center}

 Vậy phương án tối ưu phải tìm là
$$ X^*=\begin{bmatrix}
20&60&0&0\\ 
20&0&0&50\\ 
0&40&60&0\\ 
\end{bmatrix} $$ 
Chi phí tối ưu là $F(x)=390$.
\diem{1}

\end{phandiem}
\end{traloi}
\end{cauhoi}


\noindent \hrulefill

{\bf Chú ý:} Cán bộ coi thi không giải thích gì thêm.
%%%%%%%%%%%%%%%%%%%%%%%%%%%%%%%%%%%%
%Phần đóng tệp đáp án và in ra
\Closesolutionfile{goiytraloi}

%%In đáp án
\newpage
\begin{center}
\textbf{\tendhqg} \\
\textbf{ \tentruong} \\                
 -----------------------\\
\textbf{\dapan}\\
\textbf{\tenkythi, \namhoc}\\
\centerline{\textbf{Môn thi:} \textbf{\tenmonhoc}}
\end{center}

\leftline{Mã môn học: \textbf{\mamonhoc}\ \hfil Số tín chỉ:  \textbf{\sotinchi} \hfil Đề số: \textbf{\deso} \hfil }
\leftline{Dành cho sinh viên khoá: \textbf{\svkhoa} \hfil Ngành học: \textbf{\nganhhoc}\hfil }

\small
\Readsolutionfile{goiytraloi}

\begin{flushright}
\parbox[c]{8cm}{\centering
\ngaylam\\
NGƯỜI LÀM ĐÁP ÁN\\
(ký và ghi rõ họ tên)\\

\vspace*{1.5cm}
\nguoilamdapan
}
\end{flushright}
\end{document}